% ---------------------------------------------------------------
% Macros carried over verbatim from mit_stats_fundamentals.tex
% (the notes this book is built from). Do not edit here; edit the
% notes and re-copy, so the two stay in step.
% ---------------------------------------------------------------
% For the power set symbol
\usepackage[mathscr]{euscript}
\let\euscr\mathscr \let\mathscr\relax% just so we can load this and rsfs
\usepackage[scr]{rsfso}
\newcommand{\powerset}{\raisebox{.15\baselineskip}{\Large\ensuremath{\wp}}}


% Tikz settings optimized for causal graphs.
\usetikzlibrary{shapes,decorations,arrows,calc,arrows.meta,fit,positioning}
\tikzset{
    -Latex,auto,node distance =1 cm and 1 cm,semithick,
    state/.style ={ellipse, draw, minimum width = 0.7 cm},
    point/.style = {circle, draw, inner sep=0.04cm,fill,node contents={}},
    bidirected/.style={Latex-Latex,dashed},
    el/.style = {inner sep=2pt, align=left, sloped}
}


% Command to put circle around numbers in math mode
\makeatletter
\newcommand\mathcircled[1]{%
  \mathpalette\@mathcircled{#1}%
}
\newcommand\@mathcircled[2]{%
  \tikz[baseline=(math.base)] \node[draw,circle,inner sep=1pt] (math) {$\m@th#1#2$};%
}
\makeatother

% Command for strikes through arrows
\makeatletter %% <- make @ usable in command names
\newcommand*\Neg[2][0mu]{\Neginternal{#1}{\negslash}{#2}}
\newcommand*\sNeg[2][0mu]{\Neginternal{#1}{\snegslash}{#2}}
\newcommand*\ssNeg[2][0mu]{\Neginternal{#1}{\ssnegslash}{#2}}
\newcommand*\sssNeg[2][0mu]{\Neginternal{#1}{\sssnegslash}{#2}}
\newcommand*\Neginternal[3]{\mathpalette\Neg@{{#1}{#2}{#3}}}
\newcommand*\Neg@[2]{\Neg@@{#1}#2}
\newcommand*\Neg@@[4]{%
  \mathrel{\ooalign{%
    $\m@th#1#4$\cr
    \hidewidth$\m@th#3{#1}\mkern\muexpr#2*2$\hidewidth\cr
  }}%
}

\newcommand*\negslash[1]{\m@th#1\not\mathrel{\phantom{=}}}
\newcommand*\snegslash[1]{\rotatebox[origin=c]{60}{$\m@th#1-$}}
\newcommand*\ssnegslash[1]{\rotatebox[origin=c]{60}{$\m@th#1{\dabar@}\mkern-7mu{\dabar@}$}}
\newcommand*\sssnegslash[1]{\rotatebox[origin=c]{60}{$\m@th#1\dabar@$}}
\makeatother  %% <- revert @


%%%%%%%% 	FOR WIDE HATS
% \stackMath
\newcommand\reallywidehat[1]{%
\savestack{\tmpbox}{\stretchto{%
  \scaleto{%
    \scalerel*[\widthof{\ensuremath{#1}}]{\kern-.6pt\bigwedge\kern-.6pt}%
    {\rule[-\textheight/2]{1ex}{\textheight}}%WIDTH-LIMITED BIG WEDGE
  }{\textheight}% 
}{0.5ex}}%
\stackon[1pt]{#1}{\tmpbox}%
}


%%%%% COMMANDS
\DeclareMathOperator*{\E}{\mathbb{E}} % Expectation symbol
\DeclareMathOperator*{\var}{\mathbb{V}\text{ar}} % Variance
\DeclareMathOperator*{\cov}{\mathbb{C}\text{ov}} % Covariance
\DeclareMathOperator*{\corr}{\textrm{corr}} % Correlation
\DeclareMathOperator*{\se}{\textrm{se}} % Standard error
\newcommand\myeq{\stackrel{\mathclap{\normalfont\mbox{s}}}{~}} % Permits words above \sim
\newcommand{\simple}{\ensuremath{y = \beta_0 + \beta_1 x_1 + u}}
\newcommand{\multiple}{\ensuremath{y = \beta_0 + \beta_1 x_1 + \cdots +  \beta_k x_k + u}}
\newcommand{\ydev}{\ensuremath{(y_i - \bar{y})}}
\newcommand{\xdev}{\ensuremath{(x_i - \bar{x})}}
\newcommand{\xsquare}{\ensuremath{(x_i - \bar{x})^2}}
\newcommand{\vecu}{\ensuremath{\mathbf{u}}}
\newcommand{\vecut}{\ensuremath{\mathbf{u}^{\mathbf{T}}}}
\newcommand{\vecbeta}{\ensuremath{\bm{\beta}}}
\newcommand{\vecbetat}{\ensuremath{\bm{\beta}^{\mathbf{T}}}}
\newcommand{\matx}{\ensuremath{\mathbf{X}}}
\newcommand{\matxt}{\ensuremath{\mathbf{X}^{\mathbf{T}}}}
\newcommand{\mata}{\ensuremath{\mathbf{A}}}
\newcommand{\matat}{\ensuremath{\mathbf{A}^{\mathbf{T}}}}
\newcommand{\matb}{\ensuremath{\mathbf{B}}}
\newcommand{\matbt}{\ensuremath{\mathbf{B}^{\mathbf{T}}}}
\newcommand{\mathh}{\ensuremath{\mathbf{H}}}
\newcommand{\mathht}{\ensuremath{\mathbf{H}^{\mathbf{T}}}}
\newcommand{\matv}{\ensuremath{\mathbf{V}}}
\newcommand{\matvt}{\ensuremath{\mathbf{V}^{\mathbf{T}}}}
\newcommand{\maty}{\ensuremath{\mathbf{Y}}}
\newcommand{\matyt}{\ensuremath{\mathbf{Y}^{\mathbf{T}}}}
\newcommand{\matt}{\ensuremath{\mathbf{T}}}
\newcommand{\matzero}{\ensuremath{\mathbf{0}}}
\newcommand{\matzerot}{\ensuremath{\mathbf{0}^{\mathbf{T}}}}
\newcommand{\vecy}{\ensuremath{\mathbf{y}}}
\newcommand{\vecyt}{\ensuremath{\mathbf{y}^{\mathbf{T}}}}
\newcommand{\vecycap}{\ensuremath{\mathbf{Y}}}
\newcommand{\vecx}{\ensuremath{\mathbf{x}}}
\newcommand{\vecxt}{\ensuremath{\mathbf{x}^{\mathbf{T}}}}
\newcommand{\veca}{\ensuremath{\mathbf{a}}}
\newcommand{\vecat}{\ensuremath{\mathbf{a}^{\mathbf{T}}}}
\newcommand{\vecb}{\ensuremath{\mathbf{b}}}
\newcommand{\vecbt}{\ensuremath{\mathbf{b}^{\mathbf{T}}}}
\newcommand{\vecbetah}{\ensuremath{\bm{\hat{\beta}}}}
\newcommand{\vecbetatil}{\ensuremath{\bm{\tilde{\beta}}}}
\newcommand{\vecep}{\ensuremath{\bm{\epsilon}}}
\newcommand{\veceta}{\ensuremath{\bm{\eta}}}
\newcommand{\veczero}{\ensuremath{\mathbf{0}}}
\newcommand{\vectheta}{\ensuremath{{\boldsymbol \theta }}}
\newcommand{\ethn}{\ensuremath{\hat{\theta}_n}} % hat theta_n
\newcommand{\epn}{\ensuremath{\hat{p}_n}} % hat p_n
\newcommand{\xdesign}{\mathbb{X}}
\newcommand{\ident}{\mathbb{I}}
\newcommand{\Gauss}{\text{Gauss}} % Gauss text
\newcommand{\Bern}{\text{Bern}} % Bern text
\newcommand{\Pois}{\text{Pois}} % Pois text
\newcommand{\Expo}{\text{Expo}} % Expo text
\newcommand{\Bias}{\text{Bias}} % Bias text
\newcommand{\statmodel}{\ensuremath{(E, (\Pr_{\theta})_{\theta \in \Theta})}} % Stat model
\newcommand{\one}{\ensuremath{\mathds{1}}}
\renewcommand{\Pr}{\ensuremath{\text{Pr}}}
\newcommand{\card}{\ensuremath{\text{card}}}
\newcommand{\kinz}{\ensuremath{k \in \mathbb{Z}}}
\newcommand{\ninz}{\ensuremath{n \in \mathbb{Z}}}
\newcommand{\irr}{\mathbb{R}\setminus\mathbb{Q}}
\newcommand{\inz}{\ensuremath{\in \mathbb{Z}}}
\newcommand{\ps}{\ensuremath{\mathscr{P}}} % script power set
\newcommand{\dom}{\ensuremath{\text{dom}}}
\newcommand{\ran}{\ensuremath{\text{ran}}}
\newcommand{\codom}{\ensuremath{\text{codom}}}
\newcommand{\id}{\ensuremath{\text{id}}}
\newcommand{\ds}{\displaystyle}
\newcommand{\inn}{\ensuremath{\in \mathbb{N}}}
\newcommand{\xbin}{\ensuremath{X \sim \text{Bin}(n, p)}}
\DeclareMathOperator*{\sd}{\textrm{sd}} % Standard deviation
\newcommand{\intbar}{\ensuremath{\biggr\rvert}} % Evaluation bar for integrals
\renewcommand{\Pr}{\mathbb{P}}
\newcommand{\Qr}{\mathbb{Q}}
\newcommand{\real}{\ensuremath{\mathbb{R}}} % set of real numbers
\newcommand{\lv}{\ensuremath{\left \lvert}} % left bar
\newcommand{\rv}{\ensuremath{\right \rvert}} % left bar
\newcommand{\dx}{\ensuremath{\, \mathrm{d} x}} % Space between integrand and dx and non-italic d
\newcommand{\dy}{\ensuremath{\, \mathrm{d} y}} % Space between integrand and dx and non-italic d
\newcommand{\Xton}{\ensuremath{X_1, \ldots, X_n}} % Sequence of r.v.s X1 to Xn
\newcommand{\limn}{\ensuremath{\lim_{n \to \infty}}} % lim of n to infitiny
\newcommand{\iid}{\ensuremath{\stackrel{iid}{\sim}}} % iid above tilde
\newcommand{\ind}{\mathds{1}} % Indicator variables
% Convergence Commands
\newcommand{\convd}{\ensuremath{\xrightarrow [n \to \infty ]{\text{(d)}}}} % convergence in distribution
\newcommand{\convp}{\ensuremath{\xrightarrow [n \to \infty ]{\Pr}}} % convergence in prob
\newcommand{\convas}{\ensuremath{\xrightarrow [n \to \infty ]{\text{a.s.}}}} % convergence almost surely

% For longer absolute value bars (need mathtools package)
\DeclarePairedDelimiter\abs{\lvert}{\rvert}%
\DeclarePairedDelimiter\norm{\lVert}{\rVert}%

\DeclarePairedDelimiter{\oldnormaux}{\bracevert}{\bracevert}

\NewDocumentCommand{\oldnorm}{som}{%
  \IfBooleanTF{#1}
    {\oldnormaux*{#3}}
    {\IfNoValueTF{#2}
       {\oldnormaux*{\vphantom{dq}#3}}
       {\oldnormaux[#2]{#3}}%
    }%
}
